\input{preamble}

% Bon, commençons ici.
%
\begin{document}

\title{Guide bibliographique}

\maketitle

\phantomsection
\label{section-phantom}

\tableofcontents

\section{Courts articles d'introduction}
\label{section-short-introductions}


\begin{itemize}
\item Barbara Fantechi: \emph{Stacks for Everybody} \cite{fantechi_stacks}
\item Dan Edidin: \emph{What is a stack?} \cite{edidin_whatis}
\item Dan Edidin: \emph{Notes on the construction of the moduli space of
curves} \cite{edidin_notes}
\item Angelo Vistoli: \emph{Intersection theory on algebraic
stacks and on their moduli spaces}, et surtout l'appendice.
\cite{vistoli_intersection}
\end{itemize}


\section{Références classiques}
\label{section-classics}

\begin{itemize}
\item Mumford: \emph{Picard groups of moduli problems}
\cite{mumford_picard}
\begin{quote}
Mumford n'emploie jamais ici le terme « champ », mais la
notion est implicite dans l'article ; il calcule le groupe de Picard du champ
de modules des courbes elliptiques.
\end{quote}
\item Deligne, Mumford: \emph{The irreducibility of the space of curves of
given genus} \cite{DM}
\begin{quote}
Cet article influent introduit les « champs algébriques » au sens de ce que
l'on appelle aujourd'hui universellement des champs de Deligne--Mumford (champs à
diagonale représentable qui admettent des présentations étales par des schémas).
Il contient de nombreux résultats fondateurs \emph{sans démonstration}. L'article
utilise les champs pour donner deux démonstrations de l'irréductibilité de l'espace de modules des courbes de genre
$g$.
\end{quote}
\item Artin: \emph{Versal deformations and algebraic stacks}
\cite{ArtinVersal}
\begin{quote}
Cet article introduit les « champs algébriques » qui généralisent les champs de
Deligne--Mumford et que l'on appelle aujourd'hui couramment \emph{champs d'Artin} :
ce sont les champs à diagonale représentable qui admettent des présentations
lisses par des schémas. L'article donne un critère fondé sur la théorie des déformations,
connu sous le nom de critère d'Artin, qui permet de démontrer qu'un champ de
modules donné est un champ d'Artin sans exhiber explicitement de présentation.
\end{quote}
\end{itemize}




\section{Livres et notes en ligne}
\label{section-books}


\begin{itemize}
\item Laumon, Moret-Bailly: \emph{Champs Alg\'ebriques} \cite{LM-B}
\begin{quote}
Ce livre constitue actuellement la référence la plus exhaustive sur les champs
et contient de nombreux résultats fondateurs. Il suppose que le lecteur connaît
les espaces algébriques et renvoie fréquemment au livre de Knutson \cite{Kn}. Le
chapitre 12 contient une erreur concernant la fonctorialité du site lisse-étale
d'un champ algébrique. Cette erreur ne prête plus à conséquence, car Martin
Olsson l'a corrigée (voir \cite{olsson_sheaves}) et les résultats des chapitres
suivants sont justes, moyennant éventuellement de légères modifications.
\end{quote}
\item The Stacks Project Authors: \emph{Stacks Project}
\cite{stacks-project}.
\begin{quote}
Vous êtes en train de le lire !
\end{quote}
\item Anton Geraschenko:
\emph{Lecture notes for Martin Olsson's class on stacks} \cite{olsson_stacks}
\begin{quote}
Ce cours développe systématiquement la théorie des espaces
algébriques avant d'introduire les champs algébriques (définis pour la première
fois au cours 27 !). Outre les propriétés fondamentales, il traite de
l'équivalence entre être de Deligne--Mumford et avoir une diagonale non ramifiée,
du site lisse-étale d'un champ d'Artin, de la théorie des faisceaux quasi-cohérents,
du théorème de Keel--Mori, de la descente cohomologique et des gerbes (ainsi que
de leur relation avec le groupe de Brauer). Il propose également des exercices.
\end{quote}
\item
Behrend, Conrad, Edidin, Fantechi, Fulton, G\"ottsche et Kresch :
\emph{Algebraic stacks}, notes en ligne pour un livre en cours de rédaction
\cite{stacks_book}
\begin{quote}
Ce livre vise à offrir une introduction accessible aux champs, sans exiger de
prérequis approfondis, en privilégiant les exemples et les applications.
Contrairement à \cite{LM-B}, il ne suppose pas que le lecteur maîtrise la théorie
des espaces algébriques. Les champs de Deligne--Mumford y sont plutôt introduits,
les espaces algébriques constituant un cas particulier ; l'un des objectifs est de
développer une théorie suffisante pour démontrer les assertions de \cite{DM}. La
théorie générale des champs d'Artin doit être développée dans la seconde partie.
Seule une partie du livre est actuellement disponible sur le site de Kresch.
\end{quote}
\item Olsson, Martin: \emph{Algebraic spaces and stacks}, \cite{olsson_book}
\begin{quote}
Introduction vivement recommandée à la théorie des espaces algébriques et des
champs algébriques, accessible à un lecteur qui maîtrise le livre
de Hartshorne sur la géométrie algébrique.
\end{quote}
\end{itemize}


\section{Références connexes sur les fondements de la théorie des champs}
\label{section-related}


\begin{itemize}
\item Vistoli:
\emph{Notes on Grothendieck topologies, fibered categories and descent theory}
\cite{vistoli_fga}
\begin{quote}
Contient des résultats utiles sur les catégories fibrées, les champs et la
théorie de la descente pour la topologie fpqc, ainsi que des démonstrations rigoureuses.
\end{quote}
\item Knutson: \emph{Algebraic Spaces} \cite{Kn}
\begin{quote}
Ce livre, issu de sa thèse de doctorat dirigée par Michael Artin,
contient les fondements de la théorie des espaces algébriques. Le livre
\cite{LM-B} y renvoie fréquemment. Voir aussi les articles d'Artin sur les
espaces algébriques : \cite{Artin-Algebraic-Approximation},
\cite{ArtinI}, \cite{Artin-Implicit-Function},
\cite{ArtinII}, \cite{Artin-Construction-Techniques},
\cite{Artin-Algebraic-Spaces}, \cite{Artin-Theorem-Representability} et
\cite{ArtinVersal}
\end{quote}
\item Grothendieck et al., \emph{Th\'eorie des Topos et Cohomologie \'Etale des
Sch\'emas I, II, III}, également connu sous le nom de SGA 4 \cite{SGA4}
\begin{quote}
Le volume 1 contient de nombreux résultats généraux sur les univers, les sites
et les catégories fibrées. Le mot « champ » (terme français correspondant à
« stack ») apparaît dans l'exposé XVIII de Deligne.
\end{quote}
\item Jean Giraud: \emph{Cohomologie non ab\'elienne} \cite{giraud}
\begin{quote}
Le livre étudie les catégories fibrées, les champs, les torseurs et les gerbes
sur des sites généraux, mais non les champs algébriques. Par exemple, si $G$ est
un faisceau de groupes abéliens sur $X$, de même que $H^1(X, G)$ s'identifie aux
$G$-torseurs, $H^2(X, G)$ s'identifie à un ensemble convenablement défini de
$G$-gerbes. Lorsque $G$ n'est pas abélien, $H^2(X, G)$ est défini comme
l'ensemble des $G$-gerbes.
\end{quote}
\item Kelly and Street: \emph{Review of the elements of 2-categories}
\cite{kelly-street}
\begin{quote}
La catégorie des champs forme une 2-catégorie, certes d'un type simple où les
2-morphismes sont inversibles. Il s'agit d'une référence sur les 2-catégories
générales. Je ne l'ai jamais utilisée et ne puis donc en apprécier l'utilité.
Notons également que \cite{stacks-project} contient quelques notions de base sur les 2-catégories.
\end{quote}
\end{itemize}




\section{Articles de recherche}
\label{section-papers}

\noindent
On trouvera ci-dessous une liste d'articles de recherche contenant des résultats
fondamentaux sur les champs et les espaces algébriques. Les résumés visent seulement
à signaler les résultats qui contribuent à la théorie des champs ; dans bien des cas,
ces résultats sont secondaires par rapport aux objectifs principaux de l'article.
Nous répartissons les articles en catégories, certains relevant de plusieurs d'entre elles.


\subsection{Théorie des déformations et champs algébriques}
\label{subsection-deformation-theory}

\noindent
Les trois premiers articles d'Artin ne traitent pas des champs, mais ils
contiennent des résultats puissants, les deux premiers étant essentiels à
\cite{ArtinVersal}.
\begin{itemize}
\item Artin: \emph{Algebraic approximation of structures over
complete local rings} \cite{Artin-Algebraic-Approximation}
\begin{quote}
On démontre que, sous des hypothèses peu restrictives, toute déformation formelle
effective peut être approchée : si $F: (\Sch/S) \to (\textit{Sets})$
est un foncteur contravariant
localement de présentation finie, si $S$ est de type fini sur un corps ou sur un
anneau de valuation discrète excellent, si $s \in S$, et si
$\hat{\xi} \in F(\hat{\mathcal{O}}_{S, s})$ est une déformation formelle
effective, alors, pour tout $n > 0$, il existe un voisinage étale
$(S', s') \to (S, s)$ induisant un isomorphisme sur les corps résiduels, et
un élément $\xi' \in F(S')$, tels que $\xi'$ et
$\hat{\xi}$ coïncident à l'ordre $n$ (c'est-à-dire aient la même restriction dans
$F(\mathcal{O}_{S, s} / \mathfrak m^n)$).
\end{quote}
\item
Artin: \emph{Algebraization of formal moduli I} \cite{ArtinI}
\begin{quote}
On démontre que, sous des hypothèses peu restrictives, toute déformation formelle
verselle effective est algébrisable. Soit
$F: (\Sch/S) \to (\textit{Sets})$ un foncteur contravariant
localement de présentation finie, où $S$ est de type fini sur un corps ou sur un
anneau de valuation discrète excellent ; soit $s \in S$ un point localement
fermé, soit $\hat A$ une $\mathcal{O}_S$-algèbre locale noethérienne complète,
de corps résiduel $k'$, extension finie de $k(s)$,
et soit $\hat{\xi} \in F(\hat A)$ une déformation formelle verselle effective d'un
élément $\xi_0 \in F(k')$. Il existe alors un schéma $X$ de type fini sur $S$,
un point fermé $x \in X$ de corps résiduel $k(x) = k'$, et un élément
$\xi \in F(X)$, ainsi qu'un isomorphisme
$\hat{\mathcal{O}}_{X, x} \cong \hat{A}$
qui identifie les restrictions de $\xi$ et de $\hat{\xi}$ dans chaque
$F(\hat A / \mathfrak m^n)$. L'algébrisation
est unique si $\hat{\xi}$ est une déformation universelle. On en donne des
applications à la représentabilité des schémas de Hilbert et de Picard.
\end{quote}
\item Artin: \emph{Algebraization of formal moduli. II}
\cite{ArtinII}
\begin{quote}
En termes imprécis, il est démontré que, si l'on peut contracter formellement un
sous-ensemble fermé $Y' \subset X'$ au voisinage de $Y'$, alors il existe un
morphisme global $X' \to X$ qui contracte $Y$, où $X$ est un espace
algébrique.
\end{quote}
\item
Artin: \emph{Versal deformations and algebraic stacks} \cite{ArtinVersal}
\begin{quote}
Cet article décisif s'appuie sur les travaux d'Artin dans
\cite{Artin-Algebraic-Approximation} et \cite{ArtinI}. Il
introduit le critère d'Artin, qui permet d'établir l'algébricité d'un
champ en vérifiant des propriétés relevant de la théorie des déformations. Plus
précisément (mais sans prétendre à une précision parfaite), Artin construit une
présentation d'un champ $\mathcal{X}$ qui commute aux limites inductives,
au voisinage d'un point $x \in \mathcal{X}(k)$, de la manière suivante : si le
champ $\mathcal{X}$
satisfait au critère de Schlessinger (\cite{Sch}), il existe une déformation
formelle verselle
$\hat{\xi} \in \lim \mathcal{X}(\hat A / \mathfrak m^n)$ de $x$.
Si l'on suppose que les déformations formelles sont effectives (c'est-à-dire que
$\mathcal{X}(\hat{A}) \to \lim \mathcal{X}(\hat A / \mathfrak m^n)$
est bijective), on obtient une déformation formelle verselle effective
$\xi \in \mathcal{X}(\hat A)$. Les résultats de
\cite{ArtinI} fournissent alors un schéma $U$ de type fini et un
élément $\xi_U: U \to \mathcal{X}$ qui est formellement versel en un point
$u \in U$ situé au-dessus de $x$. Si l'on suppose en outre que $\mathcal{X}$
admet une théorie des déformations et des obstructions
qui satisfait certaines conditions (compatibilité avec la localisation étale
et avec la complétion, ainsi qu'une condition de constructibilité), Artin montre
à la section 4 que la versalité formelle est une condition ouverte ; quitte à
restreindre $U$, le morphisme $U \to \mathcal{X}$ est donc lisse.
Artin démontre également que tout champ admettant une présentation fppf par
un schéma admet une présentation lisse par un schéma ; on peut donc, en
particulier, former des champs quotients par des schémas en groupes plats,
séparés et de présentation finie.
\end{quote}
\item Conrad, de Jong: \emph{Approximation of Versal Deformations}
\cite{conrad-dejong}
\begin{quote}
Cet article propose une approche du théorème d'algébrisation d'Artin fondée sur
le puissant théorème de Popescu : si $A$ est un anneau noethérien et $B$ une
$A$-algèbre noethérienne, l'homomorphisme $A \to B$ est régulier si et seulement
si $B$ est une limite inductive d'$A$-algèbres lisses. On voit aisément que le
théorème de Popescu entraîne le théorème d'approximation d'Artin sur tout schéma
excellent (l'hypothèse d'excellence implique que, pour un anneau local $A$,
l'homomorphisme $A^{\text{h}} \to \hat A$ de l'hensélisé vers le complété est
régulier). L'article utilise le théorème de Popescu pour donner une
généralisation « aux groupoïdes » du théorème principal de \cite{ArtinI},
valable sur tout schéma de base excellent et en tout point $s \in S$.
En particulier, les résultats de \cite{ArtinVersal} valent sur une base
excellente arbitraire. Les auteurs étudient l'unicité locale pour la topologie
étale de l'algébrisation, ainsi que la question de savoir si le groupe des
automorphismes de l'objet agit naturellement sur l'hensélisé de
l'algébrisation.
\end{quote}
\item Jason Starr: \emph{Artin's axioms, composition, and moduli spaces}
\cite{starr_artin}
\begin{quote}
L'article établit que les axiomes d'Artin pour l'algébrisation sont compatibles
avec la composition des 1-morphismes.
\end{quote}
\item Martin Olsson: \emph{Deformation theory of representable
morphism of algebraic stacks} \cite{olsson_deformation}
\begin{quote}
Cet article étend aux morphismes représentables de champs algébriques, en termes
du complexe cotangent, les résultats usuels de la théorie des déformations des
morphismes de schémas. Ces résultats ne peuvent être considérés comme des
conséquences de la théorie générale d'Illusie, car le complexe cotangent d'un
morphisme représentable $X \to \mathcal{X}$ n'est pas
défini au moyen du complexe cotangent d'un morphisme de topos annelés (le
site lisse-étale n'étant pas fonctoriel).
\end{quote}
\end{itemize}

\subsection{Espaces de modules grossiers}
\label{subsection-coarse-moduli-spaces}

\noindent
Articles consacrés aux espaces de modules grossiers.
\begin{itemize}
\item Keel, Mori: \emph{Quotients in Groupoids} \cite{K-M}
\begin{quote}
Il semble que l'on ait longtemps tenu pour « bien connu » que les champs de
Deligne--Mumford séparés admettaient des espaces de modules grossiers. On trouve
ici une démonstration rigoureuse, quoique concise, du théorème suivant : si
$\mathcal{X}$ est un champ d'Artin localement de
type fini sur un schéma de base noethérien et si le morphisme d'inertie
$I_\mathcal{X} \to \mathcal{X}$ est fini, alors il existe un espace de
modules grossier $\phi : \mathcal{X} \to Y$
tel que $\phi$ soit séparé et que $Y$ soit un espace algébrique localement de
type fini sur $S$. La finitude de l'inertie est précisément la bonne
condition : il existe un espace de modules grossier
$\phi : \mathcal{X} \to Y$ tel que $\phi$
soit séparé si et seulement si l'inertie est finie.
\end{quote}
\item Conrad: \emph{The Keel-Mori Theorem via Stacks} \cite{conrad}
\begin{quote}
L'article de Keel et Mori \cite{K-M} est rédigé dans le langage des groupoïdes,
que certains trouvent difficile à maîtriser. Brian Conrad en donne une version
formulée dans le langage des champs, dont la démonstration est fort transparente
bien qu'elle emploie un formalisme élaboré. Conrad supprime également
l'hypothèse noethérienne.
\end{quote}
\item Rydh: \emph{Existence of quotients by finite groups and coarse moduli
spaces} \cite{rydh_quotients}
\begin{quote}
Rydh supprime l'hypothèse, imposée dans \cite{K-M} et \cite{conrad}, selon
laquelle $\mathcal{X}$
doit être de présentation finie sur une base.
\end{quote}
\item
Abramovich, Olsson, Vistoli: \emph{Tame stacks in positive characteristic}
\cite{tame}
\begin{quote}
Les auteurs définissent un \emph{champ d'Artin modéré} comme un champ d'Artin à
inertie finie tel que, si $\phi : \mathcal{X} \to Y$ est son espace de modules
grossier, le foncteur $\phi_*$ soit exact
sur les faisceaux quasi-cohérents. Ils démontrent que, pour un champ d'Artin à
inertie finie, les conditions suivantes sont équivalentes : $\mathcal{X}$ est
modéré ; les stabilisateurs de $\mathcal{X}$ sont linéairement réductifs ;
localement pour la topologie étale sur l'espace de
modules grossier, $\mathcal{X}$ est le quotient d'un schéma affine par un schéma
en groupes linéairement réductif. L'espace de modules grossier d'un champ
d'Artin modéré possède des propriétés particulièrement favorables.
Ainsi, sa formation commute à tout changement de base, tandis que
la formation de l'espace de modules grossier d'un champ d'Artin à inertie finie
ne commute en général qu'aux changements de base plats.
\end{quote}
\item Alper: \emph{Good moduli spaces for Artin stacks} \cite{alper_good}
\begin{quote}
Pour les champs d'Artin généraux dont les groupes de stabilisateurs affines sont
infinis (ces champs étant nécessairement non séparés), les espaces de modules
grossiers n'existent souvent pas. L'exemple le plus simple est
$[\mathbf{A}^1 / \mathbf{G}_m]$. On dit ici qu'un morphisme quasi-compact
$\phi : \mathcal{X} \to Y$ définit un \emph{bon espace de modules} si
$\mathcal{O}_Y \to \phi_*
\mathcal{O}_\mathcal{X}$ est un isomorphisme et si $\phi_*$ est exact sur les
faisceaux quasi-cohérents.
Cette notion généralise celle de champ d'Artin modéré de \cite{tame} et
englobe la théorie géométrique des invariants de Mumford : si $G$ est un groupe
réductif agissant linéairement sur $X \subset \mathbf{P}^n$, alors le morphisme
du champ quotient du lieu semi-stable vers le quotient GIT,
$[X^{ss}/G] \to X//G$, définit un bon espace de modules. La notion de bon
espace de modules possède de nombreuses propriétés géométriques remarquables :
(1) $\phi$ est surjectif, universellement fermé et universellement
submersif ; (2) $\phi$ identifie les points de $Y$ aux classes d'équivalence
par adhérence des points de $\mathcal{X}$ ; (3) $\phi$ est universel parmi les
morphismes vers des espaces algébriques ; (4)
les bons espaces de modules sont stables par tout changement de base ; et (5)
un fibré vectoriel sur un champ d'Artin descend au bon espace de modules si et
seulement si les représentations correspondantes sont triviales aux points
fermés.
\end{quote}
\end{itemize}


\subsection{Théorie de l'intersection}
\label{subsection-intersection-theory}

\noindent
Articles consacrés à la théorie de l'intersection sur les champs algébriques.
\begin{itemize}
\item
Vistoli: \emph{Intersection theory on algebraic stacks and on their moduli
spaces} \cite{vistoli_intersection}
\begin{quote}
Cet article établit les fondements d'une théorie de l'intersection à
coefficients rationnels pour les champs de Deligne--Mumford. Si $\mathcal{X}$
est un champ de Deligne--Mumford séparé, le groupe de Chow
$\CH_*(\mathcal{X})$ à coefficients rationnels est
défini comme le quotient du groupe abélien libre engendré par les sous-champs
fermés intègres de dimension $k$ par l'équivalence rationnelle.
On dispose d'une image réciproque plate, d'une image directe propre
et d'un homomorphisme de Gysin généralisé pour les immersions locales
régulières. Si $\phi : \mathcal{X} \to Y$ est un espace de modules
(c'est-à-dire un morphisme propre bijectif sur les
points géométriques), il existe une image directe induite
$\CH_*(\mathcal{X}) \to \CH_k(Y)$
qui est un isomorphisme.
\end{quote}
\item Edidin, Graham: \emph{Equivariant Intersection Theory}
\cite{edidin-graham}
\begin{quote}
Le but de cet article est de développer une théorie de l'intersection à
coefficients entiers pour le champ quotient $[X/G]$ associé à l'action d'un
groupe algébrique $G$ sur un espace algébrique $X$, autrement dit une théorie
de l'intersection $G$-équivariante de $X$. Définir les groupes de Chow
équivariants à l'aide des seuls cycles invariants ne donne pas une théorie
possédant de bonnes propriétés. Les auteurs généralisent plutôt la définition de
Totaro dans le cas de $BG$ et s'appuient sur le fait que, si
$V \to X$ est un fibré vectoriel, alors
$\CH_i(X) \cong \CH_i(V)$ canoniquement. Ils définissent
$\CH_i^G(X)$ comme suit :
posons $\dim(X) = n$ et $\dim(G) = g$. Pour chaque $i$, choisissons une
représentation de dimension $l$ du groupe $G$ sur un espace $V$, où $G$ agit librement dans
un ouvert $U \subset V$ dont le complémentaire est de codimension
$d > n - i$. Alors $X_G = [X \times U / G]$ est un
espace algébrique (on peut même choisir les données de sorte que ce soit un
schéma). On pose
$\CH_i^G(X) = \CH_{i + l - g}(X_G)$. Pour le champ quotient, on définit
$\CH_i( [X/G]) = \CH_{i + g}^G(X) = \CH_{i + l}(X_G)$.
En particulier, $\CH_i([X/G]) = 0$ pour $i > \dim [X/G] = n - g$,
mais ce groupe peut être non nul pour $i \ll 0$. Par exemple,
$\CH_i(B \mathbf{G}_m) = \mathbf{Z}$ pour $i \le 0$.
Les auteurs établissent que ces groupes de Chow équivariants possèdent les mêmes
propriétés fonctorielles que les groupes de Chow ordinaires. Ils montrent en
outre qu'un isomorphisme
$[X / G] \cong [Y / H]$ entraîne
$\CH_i([X/G]) = \CH_i([Y/H])$ ; la définition
ne dépend donc pas de la présentation du champ comme champ quotient.
\end{quote}
\item
Kresch: \emph{Cycle Groups for Artin Stacks} \cite{kresch_cycle}
\begin{quote}
Kresch définit les groupes de Chow des champs d'Artin arbitraires ; dans le cas
d'un champ quotient, sa définition coïncide avec celle d'Edidin et Graham dans
\cite{edidin-graham}. Pour les champs algébriques dont les groupes de
stabilisateurs sont affines, la théorie possède les propriétés usuelles.
\end{quote}
\item Behrend et Fantechi: \emph{The intrinsic normal cone}
\cite{behrend-fantechi}
\begin{quote}
Généralisant une construction due à Li et Tian, Behrend et Fantechi construisent
une classe fondamentale virtuelle pour tout champ de Deligne--Mumford.
\end{quote}
\end{itemize}


\subsection{Champs quotients}
\label{subsection-quotient-stacks}

\noindent
Les champs quotients\footnote{Dans la littérature,
\emph{champ quotient} désigne souvent un champ de la
forme $[X/G]$, où $X$ est un espace algébrique et $G$ un sous-schéma en groupes
de $\text{GL}_n$, plutôt qu'un schéma en groupes plat quelconque.}
forment une sous-classe très importante des champs d'Artin, qui comprend presque
tous les champs de modules étudiés par les géomètres algébristes. La géométrie
du champ quotient $[X/G]$ est la géométrie $G$-équivariante de $X$. Il est souvent
plus facile d'établir des propriétés pour les champs quotients, et certains
résultats ne sont connus que pour eux. Les articles suivants abordent ces questions :
quand un champ algébrique est-il un champ quotient global ? Un champ algébrique
est-il « localement » un champ quotient ?
\begin{itemize}
\item Laumon, Moret-Bailly: \cite[chapitre 6]{LM-B}
\begin{quote}
Le chapitre 6 contient plusieurs résultats sur la structure locale et globale
des champs algébriques. Il y est démontré qu'un champ algébrique $\mathcal{X}$ sur $S$
est un
champ quotient $[Y/G]$, où $Y$ est un espace algébrique (resp. un schéma, resp. un
schéma affine) et $G$ un groupe fini, si et seulement s'il existe un espace algébrique
(resp. un schéma, resp. un schéma affine) $Y'$ et un morphisme fini étale $Y'
\to \mathcal{X}$. On montre que, pour tout champ de Deligne--Mumford sur $S$ et tout
$x : \Spec(K) \to \mathcal{X}$, il existe un morphisme représentable, étale et
séparé $\phi : [X/G] \to \mathcal{X}$, où $G$ est un groupe fini
agissant sur un schéma affine sur $S$, tel
que $\Spec(K) = [X/G] \times_\mathcal{X} \Spec(K)$.
L'existence de présentations
à fibres géométriquement connexes est également étudiée en détail.
\end{quote}
\item
Edidin, Hassett, Kresch, Vistoli:
\emph{Brauer Groups and Quotient stacks} \cite{ehkv}
\begin{quote}
Ils établissent d'abord quelques faits fondamentaux (quoique peu difficiles)
concernant
les conditions sous lesquelles un champ algébrique donné (toujours supposé de type
fini sur un schéma noethérien dans cet article) est un champ quotient. Pour un champ
algébrique $\mathcal{X}$ : $\mathcal{X}$ est un champ quotient si et seulement s'il
existe un fibré vectoriel $V \to \mathcal{X}$ tel que, pour tout
point géométrique, le stabilisateur agisse fidèlement sur la fibre,
si et seulement s'il existe un fibré vectoriel $V \to \mathcal{X}$ et
un sous-champ localement fermé
$V^0 \subset V$ tel que $V^0$ soit représentable et se projette surjectivement sur
$\mathcal{X}$. Ils
établissent qu'un champ algébrique est un champ quotient s'il existe un revêtement
fini et plat par un espace algébrique. Tout champ de Deligne--Mumford lisse à
stabilisateur génériquement trivial est un champ quotient.
Ils montrent qu'une gerbe en $\mathbf{G}_m$ sur un schéma noethérien $X$
correspondant à
$\beta \in H^2(X, \mathbf{G}_m)$ est un champ quotient si et seulement si $\beta$
appartient à
l'image de l'application de Brauer $\text{Br}(X) \to \text{Br}'(X)$. Ils en déduisent
un exemple de
champ de Deligne--Mumford non séparé qui n'est pas un champ quotient.
\end{quote}
\item Totaro: \emph{The resolution property for schemes and stacks}
\cite{totaro_resolution}
\begin{quote}
Un champ possède la propriété de résolution si tout faisceau cohérent est quotient
d'un fibré vectoriel. Le premier théorème principal affirme que, si $\mathcal{X}$ est
un champ algébrique normal noethérien dont les groupes stabilisateurs aux points
fermés sont affines, alors les conditions suivantes sont équivalentes : (1)
$\mathcal{X}$ possède la propriété de résolution et
(2)
$\mathcal{X} = [Y/\text{GL}_n]$, où $Y$ est quasi-affine. Lorsque
$\mathcal{X}$ est de type fini sur
un corps, (1) et (2) équivalent aussi à : (3)
$\mathcal{X} = [\Spec(A)/G]$, où $G$ est
un schéma en groupes affine de type fini sur $k$. L'implication selon laquelle les
champs quotients possèdent la propriété de résolution a été démontrée par Thomason.
Le second théorème principal affirme que, si $\mathcal{X}$ est un champ de
Deligne--Mumford lisse sur
un corps, dont le groupe stabilisateur $I_\mathcal{X}
\to \mathcal{X}$ est fini et génériquement trivial, et dont l'espace de modules
grossier est un schéma à diagonale affine, alors
$\mathcal{X}$ possède la propriété de résolution. Un autre résultat remarquable
affirme que, si $\mathcal{X}$ est
un champ algébrique noethérien possédant la propriété de résolution, alors
$\mathcal{X}$ a
une diagonale affine si et seulement si les stabilisateurs des points fermés sont affines.
\end{quote}
\item Kresch: \emph{On the Geometry of Deligne-Mumford Stacks}
\cite{kresch_geometry}
\begin{quote}
Cet article résume des résultats généraux sur la structure des champs de
Deligne--Mumford (de type fini sur un corps) et contient plusieurs résultats
intéressants sur les champs quotients. Il y est démontré que tout champ de
Deligne--Mumford lisse, séparé et génériquement modéré, dont l'espace de modules
grossier est quasi-projectif, est un champ quotient $[Y/G]$, où $Y$ est
quasi-projectif et $G$ un groupe algébrique. Si $\mathcal{X}$ est un champ de
Deligne--Mumford dont l'espace de modules grossier
est un schéma, alors $\mathcal{X}$ est localement pour la topologie de Zariski un
champ quotient si et seulement s'il
admet un recouvrement ouvert de Zariski par des champs quotients de schémas par des
groupes finis. Si $\mathcal{X}$ est un champ de Deligne--Mumford propre sur un corps
de caractéristique 0
et d'espace de modules grossier $Y$, alors les conditions suivantes sont équivalentes :
$Y$ est projectif et $\mathcal{X}$ est un champ quotient; $Y$ est projectif et
$\mathcal{X}$ possède un faisceau générateur; enfin,
$\mathcal{X}$ admet une immersion fermée dans un champ de Deligne--Mumford lisse et
propre dont l'espace de modules grossier est projectif. Cela conduit à dire qu'un
champ de Deligne--Mumford est \emph{projectif} s'il existe une immersion fermée dans
un champ de Deligne--Mumford lisse et propre dont l'espace de modules grossier est projectif.
\end{quote}
\item Kresch, Vistoli \emph{On coverings of Deligne-Mumford stacks and
surjectivity of the Brauer map} \cite{kresch-vistoli}
\begin{quote}
Il est démontré qu'en caractéristique 0 et pour un entier $n$ fixé, les deux
assertions suivantes sont équivalentes : (1) tout champ de Deligne--Mumford lisse de
dimension $n$ est un champ quotient, et (2) le groupe de Brauer d'Azumaya coïncide
avec le groupe de Brauer cohomologique pour les schémas lisses de dimension $n$.
\end{quote}
\item Kresch: \emph{Cycle Groups for Artin Stacks} \cite{kresch_cycle}
\begin{quote}
Il est démontré que tout champ d'Artin réduit de type fini sur un corps, dont les
groupes stabilisateurs sont affines, admet une stratification par champs quotients.
\end{quote}
\item Abramovich-Vistoli:
\emph{Compactifying the space of stable maps} \cite{abramovich-vistoli}
\begin{quote}
Le lemme 2.2.3 établit que tout champ de Deligne--Mumford séparé est, localement pour
la topologie étale sur son espace de modules grossier, un champ quotient $[U/G]$, où
$U$ est affine et $G$ un groupe fini. Dans cet argument,
\cite[théorème 2.12]{olsson_homstacks} montre même que $G$ est le groupe stabilisateur.
\end{quote}
\item Abramovich, Olsson, Vistoli:
\emph{Tame stacks in positive characteristic} \cite{tame}
\begin{quote}
Cet article montre qu'un champ d'Artin modéré est, localement pour la topologie étale
sur l'espace de modules grossier, le quotient d'un schéma affine par le groupe
stabilisateur.
\end{quote}
\item Alper: \emph{On the local quotient structure of Artin stacks}
\cite{alper_quotient}
\begin{quote}
On conjecture que, pour un champ d'Artin $\mathcal{X}$ et un point fermé $x
\in \mathcal{X}$
à stabilisateur linéairement réductif, il existe un morphisme étale $[V/G_x]
\to \mathcal{X}$, où $V$ est un espace algébrique. Plusieurs éléments à l'appui de
cette conjecture sont
donnés. Un argument simple de théorie des déformations (fondé sur les idées de
\cite{tame}) montre que l'assertion est vraie dans le voisinage formel. Une
démonstration, dans le langage des champs, du théorème de la tranche étale de Luna
est présentée; elle établit que, pour les champs
$\mathcal{X} = [\Spec(A)/G]$
avec $G$ linéairement réductif, $\mathcal{X}$ est, localement pour la topologie étale
sur le quotient GIT $\Spec(A^G)$, un champ quotient par le stabilisateur.
\end{quote}
\end{itemize}

\subsection{Cohomologie}
\label{subsection-cohomology}

\noindent
Articles consacrés à la cohomologie des faisceaux sur les champs algébriques.
\begin{itemize}
\item Olsson: \emph{Sheaves on Artin stacks} \cite{olsson_sheaves}
\begin{quote}
Cet article développe la théorie des faisceaux quasi-cohérents et constructibles
et démontre les propriétés cohomologiques fondamentales. Il corrige une erreur de
\cite{LM-B} concernant la fonctorialité du site lisse-étale. Le complexe cotangent
y est construit. En outre, les théorèmes suivants y sont démontrés :
le théorème fondamental de Grothendieck pour les morphismes propres, le théorème
d'existence de Grothendieck, le théorème de connexité de Zariski et un théorème de
finitude pour les images directes propres de faisceaux cohérents et constructibles.
\end{quote}
\item Behrend: \emph{Derived $l$-adic categories for algebraic stacks}
\cite{behrend_derived}
\begin{quote}
Démontre la formule des traces de Lefschetz pour les champs algébriques.
\end{quote}
\item Behrend: \emph{Cohomology of stacks} \cite{behrend_cohomology}
\begin{quote}
Définit la cohomologie de de Rham des champs différentiables et la cohomologie
singulière des champs topologiques.
\end{quote}
\item Faltings:
\emph{Finiteness of coherent cohomology for proper fppf stacks}
\cite{faltings_finiteness}
\begin{quote}
Démontre la cohérence des images directes de faisceaux cohérents par des morphismes propres.
\end{quote}
\item Abramovich, Corti, Vistoli:
\emph{Twisted bundles and admissible covers} \cite{acv}
\begin{quote}
L'appendice contient le théorème de changement de base propre pour la cohomologie
étale des champs de Deligne--Mumford modérés.
\end{quote}
\end{itemize}


\subsection{Existence de revêtements finis par des schémas}
\label{subsection-finite-covers}

\noindent
L'existence de revêtements finis de champs de Deligne--Mumford par des schémas
est un résultat important. Dans la théorie de l'intersection sur les champs de
Deligne--Mumford, c'est un ingrédient essentiel de la définition de l'image directe
propre pour les morphismes non représentables. Plusieurs
résultats sur $\overline{\mathcal{M}}_g$ reposent sur l'existence d'un revêtement
fini par un schéma \emph{lisse}, démontrée par Looijenga. Le premier résultat dans
cette direction est peut-être \cite[théorème 6.1]{seshadri_quotients},
qui traite le cadre équivariant.
\begin{itemize}
\item Vistoli: \emph{Intersection theory on algebraic stacks and on their
moduli spaces} \cite{vistoli_intersection}
\begin{quote}
Si $\mathcal{X}$ est un champ de Deligne--Mumford muni d'un espace de modules
(c'est-à-dire d'un morphisme propre
bijectif sur les points géométriques), alors il existe un morphisme fini
$X \to \mathcal{X}$ provenant d'un schéma $X$.
\end{quote}
\item Laumon, Moret-Bailly: \cite[chapitre 16]{LM-B}
\begin{quote}
Comme application du théorème principal de Zariski, le théorème 16.6 établit que,
si $\mathcal{X}$ est un champ de Deligne--Mumford de type fini sur un schéma
noethérien, alors
il existe un morphisme fini, surjectif et génériquement étale $Z \to
\mathcal{X}$,
où $Z$ est un schéma. Le corollaire 16.6.2 montre également que tout espace
algébrique normal noethérien est isomorphe à l'espace algébrique quotient $X'/G$
d'un schéma normal $X'$ par l'action d'un groupe fini $G$.
\end{quote}
\item Edidin, Hassett, Kresch, Vistoli:
\emph{Brauer Groups and Quotient stacks} \cite{ehkv}
\begin{quote}
Le théorème 2.7 affirme que,
si $\mathcal{X}$ est un champ algébrique de type fini sur un
schéma de base noethérien $S$, alors la diagonale
$\mathcal{X} \to \mathcal{X} \times_S \mathcal{X}$ est
quasi-finie si et seulement s'il existe un morphisme fini surjectif
$X \to \mathcal{X}$ provenant d'un schéma $X$.
\end{quote}
\item Kresch, Vistoli: \emph{On coverings of Deligne-Mumford stacks and
surjectivity of the Brauer map} \cite{kresch-vistoli}
\begin{quote}
Il est démontré que tout champ de Deligne--Mumford lisse et séparé, de type fini
sur un corps et dont l'espace de modules grossier est quasi-projectif, admet un
revêtement fini et plat par un schéma lisse quasi-projectif.
\end{quote}
\item Olsson: \emph{On proper coverings of Artin stacks} \cite{olsson_proper}
\begin{quote}
Il est démontré que, si $\mathcal{X}$ est un champ d'Artin séparé
et de type fini sur $S$, alors
il existe un morphisme propre surjectif $X \to \mathcal{X}$ provenant d'un schéma
$X$ quasi-projectif sur $S$. Comme application, Olsson démontre la cohérence et
la constructibilité des faisceaux images directes par les morphismes propres. Il
en déduit également le théorème d'existence de Grothendieck pour les champs
d'Artin propres.
\end{quote}
\item Rydh: \emph{Noetherian approximation of algebraic spaces and stacks}
\cite{rydh_approx}
\begin{quote}
Le théorème B de cet article s'énonce comme suit.
Soit $X$ un champ algébrique quasi-compact à diagonale quasi-finie et séparée
(resp.\ un champ de Deligne--Mumford quasi-compact à diagonale quasi-compacte
et séparée). Il existe alors un schéma $Z$ et un morphisme $Z \to X$ fini,
de présentation finie et surjectif, qui est plat (resp.\ étale) au-dessus d'un
sous-champ ouvert dense quasi-compact $U \subset X$.
\end{quote}
\end{itemize}


\subsection{Rigidification}
\label{subsection-rigidification}

\noindent
La rigidification est un procédé qui retire un sous-groupe plat de l'inertie.
Par exemple, si $X$ est une variété projective, le morphisme du champ de Picard
vers le schéma de Picard est la rigidification par le groupe d'automorphismes
$\mathbf{G}_m$.
\begin{itemize}
\item Abramovich, Corti, Vistoli :
\emph{Twisted bundles and admissible covers} \cite{acv}
\begin{quote}
Soient $\mathcal{X}$ un champ algébrique sur $S$ et $H$ un schéma en groupes
plat, séparé et de présentation finie sur $S$. Supposons que, pour tout objet
$\xi \in \mathcal{X}(T)$, il existe un plongement
$H(T) \hookrightarrow \text{Aut}_{\mathcal{X}(T)}(\xi)$ compatible aux
changements de base, au sens où, pour toute flèche
$\phi : \xi \rightarrow \xi'$ au-dessus de $f: T \rightarrow T'$ et tout
$g \in H(T')$, on a $g \circ \phi = \phi \circ f^*g$.
Il existe alors un champ algébrique $\mathcal{X}/H$ et un morphisme
$\rho : \mathcal{X} \rightarrow \mathcal{X}/H$ qui est une gerbe fppf et tel
que, pour tout $\xi \in \mathcal{X}(T)$, le morphisme
$\text{Aut}_{\mathcal{X}(T)} (\xi)
\rightarrow \text{Aut}_{\mathcal{X}/H (T)} (\xi) $
est surjectif de noyau $H(T)$.
\end{quote}
\item Romagny : \emph{Group actions on stacks and applications}
\cite{romagny_actions}
\begin{quote}
Étudie le comportement des actions de groupes vis-à-vis des rigidifications.
\end{quote}
\item Abramovich, Graber, Vistoli :
\emph{Gromov-Witten theory for Deligne-Mumford stacks} \cite{agv}
\begin{quote}
L'appendice donne une synthèse de la rigidification au sens de \cite{acv},
ainsi que deux interprétations alternatives. Cet article contient aussi des
constructions permettant de recoller des champs algébriques le long de
sous-champs fermés et de prendre des racines de fibrés en droites.
\end{quote}
\item
Abramovich, Olsson, Vistoli : \emph{Tame stacks in positive characteristic}
(\cite{tame})
\begin{quote}
L'appendice traite la situation plus complexe où le sous-champ en groupes plat
de l'inertie $H \subset I_\mathcal{X}$ est normal, mais pas nécessairement
central.
\end{quote}
\end{itemize}

\subsection{Courbes stacky}
\label{subsection-stacky-curves}

\noindent
Articles consacrés aux courbes stacky.
\begin{itemize}
\item Abramovich, Vistoli : \emph{Compactifying the space of stable maps}
\cite{abramovich-vistoli}
\begin{quote}
Cet article introduit les \emph{courbes tordues}. L'espace de modules des
applications stables de courbes stables dans un champ algébrique n'est en
général pas compact. En utilisant des applications définies sur des courbes
tordues, les auteurs construisent un champ de modules qui est propre lorsque
le but est un champ de Deligne--Mumford modéré dont l'espace de modules
grossier est projectif.
\end{quote}
\item Behrend, Noohi : \emph{Uniformization of Deligne-Mumford curves}
\cite{behrend-noohi}
\begin{quote}
Démontre un théorème d'uniformisation des courbes analytiques de
Deligne--Mumford.
\end{quote}
\end{itemize}

\subsection{Champs de Hilbert, de Quot, de Hom et de variétés branchées}
\label{subsection-hilbert-quot-hom}

\noindent
Articles consacrés aux schémas de Hilbert et aux constructions analogues.
\begin{itemize}
\item Vistoli : \emph{The Hilbert stack and the theory of moduli of families}
\cite{vistoli_hilbert}
\begin{quote}
Si $\mathcal{X}$ est un champ algébrique séparé et localement de type fini sur
un espace algébrique $S$ localement noethérien et localement séparé, Vistoli
définit le champ de Hilbert $\mathcal{H}\text{ilb}(\mathcal{X} / S)$ qui
paramètre les morphismes finis et non ramifiés provenant de schémas propres.
Il est affirmé sans démonstration que
$\mathcal{H}\text{ilb}(\mathcal{X} / S)$ est un champ algébrique. On en déduit
que, pour $\mathcal{X}$ comme ci-dessus, le champ de Hom
$\mathcal{H} \text{om}_S(T, \mathcal{X})$ est un champ algébrique si $T$ est
propre et plat sur $S$.
\end{quote}
\item Olsson, Starr : \emph{Quot functors for Deligne-Mumford stacks}
\cite{olsson-starr}
\begin{quote}
Si $\mathcal{X}$ est un champ de Deligne--Mumford séparé et localement de
présentation finie sur un espace algébrique $S$, et si $\mathcal{F}$ est un
$\mathcal{O}_\mathcal{X}$-module localement de présentation finie, le foncteur
Quot $\text{Quot}(\mathcal{F} / \mathcal{X} / S)$ est représenté par un espace
algébrique séparé et localement de présentation finie sur $S$. Cet article
définit aussi les faisceaux générateurs et démontre l'existence d'un tel
faisceau pour les champs de Deligne--Mumford modérés et séparés qui sont des
champs quotients globaux d'un schéma par un groupe fini.
\end{quote}
\item Olsson : \emph{Hom-stacks and Restrictions of Scalars}
\cite{olsson_homstacks}
\begin{quote}
Supposons que $\mathcal{X}$ et $\mathcal{Y}$ soient des champs d'Artin
localement de présentation finie sur un espace algébrique $S$, à diagonale
finie, que $\mathcal{X}$ soit propre et plat sur $S$, et que, localement pour
la topologie fppf sur $S$, $\mathcal{X}$ admette un revêtement fini, plat et de
présentation finie par un espace algébrique (par exemple si $\mathcal{X}$ est
de Deligne--Mumford ou est un champ d'Artin modéré). Alors
$\Hom_S(\mathcal{X}, \mathcal{Y})$ est un champ d'Artin localement de
présentation finie sur $S$.
\end{quote}
\item Alexeev et Knutson : \emph{Complete moduli spaces of branchvarieties}
(\cite{alexeev-knutson})
\begin{quote}
Ils définissent une variété branchée de $\mathbf{P}^n$ comme un morphisme fini
$X \rightarrow \mathbf{P}^n$ provenant d'un schéma \emph{réduit} $X$. Ils
démontrent que le champ de modules des variétés branchées de polynôme de
Hilbert fixé et de degrés totaux fixés pour les composantes de dimension $i$
est un champ d'Artin propre à stabilisateur fini. Ils comparent le champ des
variétés branchées au schéma de Hilbert, au schéma de Chow et à l'espace de
modules des applications stables.
\end{quote}
\item Lieblich : \emph{Remarks on the stack of coherent algebras}
\cite{lieblich_remarks}
\begin{quote}
Cet article construit une généralisation du champ des variétés branchées
d'Alexeev et Knutson au-dessus d'un schéma $Y$, en le réalisant comme un champ
d'algèbres sur le faisceau structural de $Y$. Il donne des démonstrations de
l'existence des espaces $\text{Quot}$ et $\Hom$.
\end{quote}
\item Starr : \emph{Artin's axioms, composition, and moduli spaces}
\cite{starr_artin}
\begin{quote}
Comme application du résultat principal, l'article fournit une généralisation
commune du champ de Hilbert de Vistoli \cite{vistoli_hilbert} et du champ des
variétés branchées d'Alexeev et Knutson \cite{alexeev-knutson}. Si
$\mathcal{X}$ est un champ algébrique localement de type fini sur un schéma
excellent $S$, à diagonale finie, alors le champ $\mathcal{H}$ qui paramètre
les morphismes $g: T \rightarrow \mathcal{X}$ provenant d'un espace algébrique
propre $T$ muni d'un fibré en droites $g$-ample $L$ est un champ d'Artin
localement de type fini sur $S$.
\end{quote}
\item Lundkvist et Skjelnes :
\emph{Non-effective deformations of Grothendieck's Hilbert functor}
\cite{lundkvist-skjelnes}
\begin{quote}
Montre que le foncteur de Hilbert d'un schéma non séparé n'est pas représenté,
car il existe des déformations non effectives.
\end{quote}
\item Halpern-Leistner et Preygel :
\emph{Mapping stacks and categorical notions of properness}
\cite{HL-P}
\begin{quote}
Cet article démontre que le champ de Hom est algébrique sous des hypothèses sur
la source et le but qui sont plus générales que celles de l'article d'Olsson,
ou du moins différentes.
\end{quote}
\end{itemize}



\subsection{Champs toriques}
\label{subsection-toric}

\noindent
Les champs toriques fournissent une vaste classe d'exemples et un banc d'essai
naturel pour les conjectures, grâce au dictionnaire entre la géométrie d'un
champ torique et la combinatoire de son éventail stacky, tout comme les
variétés toriques fournissent des exemples et des contre-exemples en théorie
des schémas.
\begin{itemize}
\item Borisov, Chen et Smith : \emph{The orbifold Chow ring of toric
Deligne-Mumford stacks} \cite{bcs}
\begin{quote}
Inspiré par la construction de Cox pour les variétés toriques, cet article
définit les champs toriques de Deligne--Mumford lisses comme des champs
quotients explicites associés à un objet combinatoire appelé
\emph{éventail stacky}.
\end{quote}
\item Iwanari : \emph{The category of toric stacks} \cite{iwanari_toric}
\begin{quote}
Cet article définit un \emph{triplet torique} comme un champ de
Deligne--Mumford lisse $\mathcal{X}$ muni d'une immersion ouverte
$\mathbf{G}_m \hookrightarrow \mathcal{X}$ d'image dense (de sorte que
$\mathcal{X}$ est un orbifold) et d'une action
$\mathcal{X} \times \mathbf{G}_m \rightarrow \mathcal{X}$. Il est démontré
qu'il existe une équivalence entre la 2-catégorie des triplets toriques et la
1-catégorie des éventails stacky. L'article étudie la relation entre les
triplets toriques et la définition des champs toriques de Deligne--Mumford
lisses donnée dans \cite{bcs}.
\end{quote}
\item Iwanari : \emph{Integral Chow rings for toric stacks}
\cite{iwanari_chow}
\begin{quote}
Généralise les $\Delta$-collections de Cox pour les variétés toriques aux
orbifolds toriques.
\end{quote}
\item Perroni : \emph{A note on toric Deligne-Mumford stacks}
\cite{perroni}
\begin{quote}
Généralise les $\Delta$-collections de Cox et l'article d'Iwanari
\cite{iwanari_chow} aux champs toriques de Deligne--Mumford lisses généraux.
\end{quote}
\item Fantechi, Mann et Nironi : \emph{Smooth toric DM stacks}
\cite{fmn}
\begin{quote}
Cet article définit un champ torique de Deligne--Mumford lisse comme un champ
de Deligne--Mumford lisse $\mathcal{X}$ muni de l'action d'un tore de
Deligne--Mumford $\mathcal{T}$ (c'est-à-dire d'un champ de Picard isomorphe à
$T \times BG$, où $G$ est fini) qui possède une orbite ouverte dense isomorphe
à $\mathcal{T}$. Les auteurs en donnent une « description ascendante » et
démontrent une équivalence entre les champs toriques de Deligne--Mumford
lisses et les éventails stacky.
\end{quote}
\item Geraschenko et Satriano : \emph{Toric Stacks I and II}
\cite{gs_toric1} et \cite{gs_toric2}
\begin{quote}
Ces articles définissent un champ torique comme le champ quotient d'une
variété torique par un sous-groupe de son tore. Un champ torique génériquement
stacky est défini comme un sous-champ invariant par le tore d'un champ
torique. Cette définition englobe et étend les définitions antérieures des
champs toriques. Le premier article établit un dictionnaire entre la
combinatoire des éventails stacky et la géométrie des champs
correspondants. Il donne aussi une interprétation modulaire des champs
toriques lisses, généralisant celle de \cite{perroni}. Le second article
démontre une caractérisation intrinsèque des champs toriques.
\end{quote}
\end{itemize}



\subsection{Théorème des fonctions formelles et théorème d'existence de Grothendieck}
\label{subsection-theorem-formal-functions}

\noindent
Ces articles donnent des généralisations du théorème des fonctions formelles
\cite[III.4.1.5]{EGA} (parfois appelé théorème fondamental de Grothendieck
pour les morphismes propres) et du théorème d'existence de Grothendieck
\cite[III.5.1.4]{EGA}.
\begin{itemize}
\item Knutson : \emph{Algebraic spaces} \cite[chapitre V]{Kn}
\begin{quote}
Généralise ces théorèmes aux espaces algébriques.
\end{quote}
\item Abramovich--Vistoli : \emph{Compactifying the space of stable maps}
\cite[A.1.1]{abramovich-vistoli}
\begin{quote}
Généralise ces théorèmes aux champs de Deligne--Mumford modérés.
\end{quote}
\item Olsson et Starr : \emph{Quot functors for Deligne-Mumford stacks}
\cite{olsson-starr}
\begin{quote}
Généralise ces théorèmes aux champs de Deligne--Mumford séparés.
\end{quote}
\item Olsson : \emph{On proper coverings of Artin stacks}
\cite{olsson_proper}
\begin{quote}
En fournit une généralisation aux champs d'Artin propres.
\end{quote}
\item Conrad : \emph{Formal GAGA on Artin stacks} \cite{conrad_gaga}
\begin{quote}
En fournit une généralisation aux champs d'Artin propres et démontre un
théorème de GAGA formel.
\end{quote}
\item Olsson : \emph{Sheaves on Artin stacks} \cite{olsson_sheaves}
\begin{quote}
Donne une autre démonstration de la généralisation aux champs d'Artin propres.
\end{quote}
\end{itemize}


\subsection{Actions de groupes sur les champs}
\label{subsection-group-actions}

\noindent
Les actions de groupes sur les champs algébriques apparaissent naturellement.
Par exemple, le groupe symétrique $S_n$ agit sur
$\overline{\mathcal{M}}_{g, n}$ et, lorsqu'un groupe $G$ agit sur un schéma
$X$, le normalisateur de $G$ dans $\text{Aut}(X)$ agit sur $[X/G]$. En outre,
les actions de tores sur les champs apparaissent souvent en théorie de
Gromov--Witten.
\begin{itemize}
\item Romagny : \emph{Group actions on stacks and applications}
\cite{romagny_actions}
\begin{quote}
Cet article précise ce que signifie l'action d'un groupe sur un champ
algébrique et démontre l'existence de points fixes ainsi que celle de
quotients pour les actions de schémas en groupes sur les champs algébriques.
Voir aussi la note antérieure de Romagny \cite{romagny_notes}.
\end{quote}
\end{itemize}

\subsection{Prise de racines de fibrés en droites}
\label{subsection-root-stacks}

\noindent
Cette construction utile a été découverte indépendamment par Cadman et par
Abramovich, Graber et Vistoli. Étant donné un schéma $X$ muni d'un diviseur de
Cartier effectif $D$, le champ des racines $r$-ièmes est un champ d'Artin
ramifié au-dessus de $X$ le long de $D$, avec un stabilisateur $\mu_r$ au-dessus
de $D$, et il est schématique hors de $D$.
\begin{itemize}
\item Charles Cadman
\emph{Using Stacks to Impose Tangency Conditions on Curves}
\cite{cadman}
\item Abramovich, Graber, Vistoli : \emph{Gromov-Witten theory for
Deligne-Mumford stacks} \cite{agv}
\end{itemize}

\subsection{Autres articles}
\label{subsection-other}

\noindent
Un assortiment d'autres articles.
\begin{itemize}
\item Lieblich : \emph{Moduli of twisted sheaves} \cite{lieblich_twisted}
\begin{quote}
Cet article contient une synthèse sur les gerbes et les faisceaux tordus.
Si $\mathcal{X} \rightarrow X$ est une $\mu_n$-gerbe, où $X$ est une surface
relative projective à fibres géométriques lisses et connexes, il est démontré
que le champ des faisceaux $\mathcal{X}$-tordus semi-stables est un champ
d'Artin localement de présentation finie sur $S$. L'article développe aussi
la théorie des points associés et de la pureté des faisceaux sur les champs
d'Artin.
\end{quote}
\item Lieblich, Osserman :
\emph{Functorial reconstruction theorem for stacks}
\cite{lieblich-osserman}
\begin{quote}
Démontre plusieurs résultats surprenants et intéressants sur les conditions
permettant de reconstruire un champ algébrique à partir du foncteur qui lui
est associé.
\end{quote}
\item David Rydh :
\emph{Noetherian approximation of algebraic spaces and stacks}
\cite{rydh_approx}
\begin{quote}
Cet article montre que tout champ algébrique quasi-compact à diagonale
quasi-finie peut être obtenu par approximation par un champ noethérien. Cela
permet notamment de supprimer l'hypothèse noethérienne dans des résultats de
Chevalley, Serre, Zariski et Chow.
\end{quote}
\end{itemize}



\section{Les champs dans d'autres domaines}
\label{section-stacks-areas}

\begin{itemize}
\item Behrend et Noohi : \emph{Uniformization of Deligne-Mumford curves}
\cite{behrend-noohi}
\begin{quote}
Donne une vue d'ensemble et une comparaison des champs topologiques,
analytiques et algébriques.
\end{quote}
\item Behrang Noohi : \emph{Foundations of topological stacks I} \cite{noohi}
\item David Metzler : \emph{Topological and smooth stacks} \cite{metzler}
\end{itemize}


\section{Champs supérieurs}
\label{section-higher-stacks}

\begin{itemize}
\item Lurie : \emph{Higher topos theory} \cite{lurie_topos}
\item Lurie : \emph{Derived Algebraic Geometry I - V}
\cite{dag1}, \cite{dag2}, \cite{dag3}, \cite{dag4}, \cite{dag5}
\item To\"en : \emph{Higher and derived stacks: a global overview}
\cite{toen_higher}
\item To\"en et Vezzosi : \emph{Homotopical algebraic geometry I, II}
\cite{hag1}, \cite{hag2}
\end{itemize}

\begin{multicols}{2}[\section{Autres chapitres}]
\noindent
Préliminaires
\begin{enumerate}
\item \hyperref[introduction-section-phantom]{Introduction}
\item \hyperref[conventions-section-phantom]{Conventions}
\item \hyperref[sets-section-phantom]{Théorie des ensembles}
\item \hyperref[categories-section-phantom]{Catégories}
\item \hyperref[topology-section-phantom]{Topologie}
\item \hyperref[sheaves-section-phantom]{Faisceaux sur les espaces}
\item \hyperref[sites-section-phantom]{Sites et faisceaux}
\item \hyperref[stacks-section-phantom]{Champs}
\item \hyperref[fields-section-phantom]{Corps}
\item \hyperref[algebra-section-phantom]{Algèbre commutative}
\item \hyperref[brauer-section-phantom]{Groupes de Brauer}
\item \hyperref[homology-section-phantom]{Algèbre homologique}
\item \hyperref[derived-section-phantom]{Catégories dérivées}
\item \hyperref[simplicial-section-phantom]{Méthodes simpliciales}
\item \hyperref[more-algebra-section-phantom]{Compléments d'algèbre}
\item \hyperref[smoothing-section-phantom]{Lissage des morphismes d'anneaux}
\item \hyperref[modules-section-phantom]{Faisceaux de modules}
\item \hyperref[sites-modules-section-phantom]{Modules sur les sites}
\item \hyperref[injectives-section-phantom]{Objets injectifs}
\item \hyperref[cohomology-section-phantom]{Cohomologie des faisceaux}
\item \hyperref[sites-cohomology-section-phantom]{Cohomologie sur les sites}
\item \hyperref[dga-section-phantom]{Algèbre différentielle graduée}
\item \hyperref[dpa-section-phantom]{Algèbre à puissances divisées}
\item \hyperref[sdga-section-phantom]{Faisceaux différentiels gradués}
\item \hyperref[hypercovering-section-phantom]{Hyperrecouvrements}
\end{enumerate}
Schémas
\begin{enumerate}
\setcounter{enumi}{25}
\item \hyperref[schemes-section-phantom]{Schémas}
\item \hyperref[constructions-section-phantom]{Constructions de schémas}
\item \hyperref[properties-section-phantom]{Propriétés des schémas}
\item \hyperref[morphisms-section-phantom]{Morphismes de schémas}
\item \hyperref[coherent-section-phantom]{Cohomologie des schémas}
\item \hyperref[divisors-section-phantom]{Diviseurs}
\item \hyperref[limits-section-phantom]{Limites de schémas}
\item \hyperref[varieties-section-phantom]{Variétés}
\item \hyperref[topologies-section-phantom]{Topologies sur les schémas}
\item \hyperref[descent-section-phantom]{Descente}
\item \hyperref[perfect-section-phantom]{Catégories dérivées de schémas}
\item \hyperref[more-morphisms-section-phantom]{Compléments sur les morphismes}
\item \hyperref[flat-section-phantom]{Compléments sur la platitude}
\item \hyperref[groupoids-section-phantom]{Schémas en groupoïdes}
\item \hyperref[more-groupoids-section-phantom]{Compléments sur les schémas en groupoïdes}
\item \hyperref[etale-section-phantom]{Morphismes étales de schémas}
\end{enumerate}
Thèmes de théorie des schémas
\begin{enumerate}
\setcounter{enumi}{41}
\item \hyperref[chow-section-phantom]{Homologie de Chow}
\item \hyperref[intersection-section-phantom]{Théorie de l'intersection}
\item \hyperref[pic-section-phantom]{Schémas de Picard des courbes}
\item \hyperref[weil-section-phantom]{Théories cohomologiques de Weil}
\item \hyperref[adequate-section-phantom]{Modules adéquats}
\item \hyperref[dualizing-section-phantom]{Complexes dualisants}
\item \hyperref[duality-section-phantom]{Dualité pour les schémas}
\item \hyperref[discriminant-section-phantom]{Discriminants et différentes}
\item \hyperref[derham-section-phantom]{Cohomologie de de Rham}
\item \hyperref[local-cohomology-section-phantom]{Cohomologie locale}
\item \hyperref[algebraization-section-phantom]{Géométrie algébrique et formelle}
\item \hyperref[curves-section-phantom]{Courbes algébriques}
\item \hyperref[resolve-section-phantom]{Résolution des surfaces}
\item \hyperref[models-section-phantom]{Réduction semi-stable}
\item \hyperref[functors-section-phantom]{Foncteurs et morphismes}
\item \hyperref[equiv-section-phantom]{Catégories dérivées de variétés}
\item \hyperref[pione-section-phantom]{Groupes fondamentaux de schémas}
\item \hyperref[etale-cohomology-section-phantom]{Cohomologie étale}
\item \hyperref[crystalline-section-phantom]{Cohomologie cristalline}
\item \hyperref[proetale-section-phantom]{Cohomologie pro-étale}
\item \hyperref[relative-cycles-section-phantom]{Cycles relatifs}
\item \hyperref[more-etale-section-phantom]{Compléments de cohomologie étale}
\item \hyperref[trace-section-phantom]{La formule des traces}
\end{enumerate}
Espaces algébriques
\begin{enumerate}
\setcounter{enumi}{64}
\item \hyperref[spaces-section-phantom]{Espaces algébriques}
\item \hyperref[spaces-properties-section-phantom]{Propriétés des espaces algébriques}
\item \hyperref[spaces-morphisms-section-phantom]{Morphismes d'espaces algébriques}
\item \hyperref[decent-spaces-section-phantom]{Espaces algébriques décents}
\item \hyperref[spaces-cohomology-section-phantom]{Cohomologie des espaces algébriques}
\item \hyperref[spaces-limits-section-phantom]{Limites d'espaces algébriques}
\item \hyperref[spaces-divisors-section-phantom]{Diviseurs sur les espaces algébriques}
\item \hyperref[spaces-over-fields-section-phantom]{Espaces algébriques sur des corps}
\item \hyperref[spaces-topologies-section-phantom]{Topologies sur les espaces algébriques}
\item \hyperref[spaces-descent-section-phantom]{Descente et espaces algébriques}
\item \hyperref[spaces-perfect-section-phantom]{Catégories dérivées d'espaces}
\item \hyperref[spaces-more-morphisms-section-phantom]{Compléments sur les morphismes d'espaces}
\item \hyperref[spaces-flat-section-phantom]{Platitude sur les espaces algébriques}
\item \hyperref[spaces-groupoids-section-phantom]{Groupoïdes dans les espaces algébriques}
\item \hyperref[spaces-more-groupoids-section-phantom]{Compléments sur les groupoïdes dans les espaces}
\item \hyperref[bootstrap-section-phantom]{Amorçage}
\item \hyperref[spaces-pushouts-section-phantom]{Sommes amalgamées d'espaces algébriques}
\end{enumerate}
Thèmes de géométrie
\begin{enumerate}
\setcounter{enumi}{81}
\item \hyperref[spaces-chow-section-phantom]{Groupes de Chow des espaces}
\item \hyperref[groupoids-quotients-section-phantom]{Quotients de groupoïdes}
\item \hyperref[spaces-more-cohomology-section-phantom]{Compléments sur la cohomologie des espaces}
\item \hyperref[spaces-simplicial-section-phantom]{Espaces simpliciaux}
\item \hyperref[spaces-duality-section-phantom]{Dualité pour les espaces}
\item \hyperref[formal-spaces-section-phantom]{Espaces algébriques formels}
\item \hyperref[restricted-section-phantom]{Algébrisation des espaces formels}
\item \hyperref[spaces-resolve-section-phantom]{Résolution des surfaces revisitée}
\end{enumerate}
Théorie des déformations
\begin{enumerate}
\setcounter{enumi}{89}
\item \hyperref[formal-defos-section-phantom]{Théorie formelle des déformations}
\item \hyperref[defos-section-phantom]{Théorie des déformations}
\item \hyperref[cotangent-section-phantom]{Le complexe cotangent}
\item \hyperref[examples-defos-section-phantom]{Problèmes de déformation}
\end{enumerate}
Champs algébriques
\begin{enumerate}
\setcounter{enumi}{93}
\item \hyperref[algebraic-section-phantom]{Champs algébriques}
\item \hyperref[examples-stacks-section-phantom]{Exemples de champs}
\item \hyperref[stacks-sheaves-section-phantom]{Faisceaux sur les champs algébriques}
\item \hyperref[criteria-section-phantom]{Critères de représentabilité}
\item \hyperref[artin-section-phantom]{Axiomes d'Artin}
\item \hyperref[quot-section-phantom]{Espaces de Quot et de Hilbert}
\item \hyperref[stacks-properties-section-phantom]{Propriétés des champs algébriques}
\item \hyperref[stacks-morphisms-section-phantom]{Morphismes de champs algébriques}
\item \hyperref[stacks-limits-section-phantom]{Limites de champs algébriques}
\item \hyperref[stacks-cohomology-section-phantom]{Cohomologie des champs algébriques}
\item \hyperref[stacks-perfect-section-phantom]{Catégories dérivées de champs}
\item \hyperref[stacks-introduction-section-phantom]{Introduction aux champs algébriques}
\item \hyperref[stacks-more-morphisms-section-phantom]{Compléments sur les morphismes de champs}
\item \hyperref[stacks-geometry-section-phantom]{La géométrie des champs}
\end{enumerate}
Thèmes de théorie des espaces de modules
\begin{enumerate}
\setcounter{enumi}{107}
\item \hyperref[moduli-section-phantom]{Champs de modules}
\item \hyperref[moduli-curves-section-phantom]{Espaces de modules de courbes}
\end{enumerate}
Divers
\begin{enumerate}
\setcounter{enumi}{109}
\item \hyperref[examples-section-phantom]{Exemples}
\item \hyperref[exercises-section-phantom]{Exercices}
\item \hyperref[guide-section-phantom]{Guide bibliographique}
\item \hyperref[desirables-section-phantom]{Desiderata}
\item \hyperref[coding-section-phantom]{Style de programmation}
\item \hyperref[obsolete-section-phantom]{Obsolète}
\item \hyperref[fdl-section-phantom]{Licence de documentation libre GNU}
\item \hyperref[index-section-phantom]{Index généré automatiquement}
\end{enumerate}
\end{multicols}



\bibliography{my}
\bibliographystyle{amsalpha}

\end{document}
